\section{算术编码：从符号码到序列区间}
\label{sec:07-Information-Theory/04-Source-Coding/03}

三个符号，\(P(A)=0.5\)、\(P(B)=0.3\)、\(P(C)=0.2\)。理想码长分别是 \(1\)、\(1.74\)、\(2.32\) 比特，而 Huffman 只能给出 \(1\)、\(2\)、\(2\)。\(B\) 被硬塞了 \(0.26\) 个多余比特，\(C\) 被硬塞了 \(0.32\) 个，两笔零头一个都退不掉。

短消息里这些零头会互相抵消，看不出问题。发一千个符号就藏不住了：熵给出的理想总长约 \(1485\) 比特，Huffman 实际用掉 \(1500\)。十五个比特不算灾难，但它是系统性的，序列越长越明显，而且靠调整码表消不掉——码表已经最优。

墙在规则里。只要坚持``每个符号配一个整数长度的码字''，每个符号就各自承担一份取整损失，损失之和就是那笔多余开销。绑成块只是把损失摊薄，字母表却按指数膨胀，追到一定精度就追不动了。

\subsection{整段序列对应一个区间，不再对应一串码字}

换个立足点。不发码字，发一个数。这个数落在 \([0,1)\) 的某个位置，位置本身就唯一地标定了整条消息。

做法是逐符号切割区间。先按累积概率把 \([0,1)\) 分成三段，宽度分别是各符号的概率。读到第一个符号，就把当前区间换成它对应的那一段；读到第二个符号，把这一段内部按同样的比例再切一次，取相应的子段。如此进行下去，区间越缩越窄。

\begin{center}
\begin{tikzpicture}[x=1cm,y=1cm,font=\small,>=stealth]
  % 第一步
  \fill[black!6] (0,2.2) rectangle (9,2.8);
  \fill[blue!25] (4.5,2.2) rectangle (7.2,2.8);
  \draw[black!50] (0,2.2) rectangle (9,2.8);
  \draw[black!50] (4.5,2.2) -- (4.5,2.8);
  \draw[black!50] (7.2,2.2) -- (7.2,2.8);
  \node[font=\footnotesize] at (2.25,2.5) {\(A\)};
  \node[font=\footnotesize] at (5.85,2.5) {\(B\)};
  \node[font=\footnotesize] at (8.1,2.5) {\(C\)};
  \node[left,font=\footnotesize] at (-0.12,2.5) {读入 \(B\)};
  \node[right,font=\footnotesize,black!65] at (9.15,2.5) {\([0.5,\,0.8)\)，宽 \(0.3\)};
  % 第二步
  \fill[black!6] (0,1.0) rectangle (9,1.6);
  \fill[blue!25] (0,1.0) rectangle (4.5,1.6);
  \draw[black!50] (0,1.0) rectangle (9,1.6);
  \draw[black!50] (4.5,1.0) -- (4.5,1.6);
  \draw[black!50] (7.2,1.0) -- (7.2,1.6);
  \node[font=\footnotesize] at (2.25,1.3) {\(A\)};
  \node[font=\footnotesize] at (5.85,1.3) {\(B\)};
  \node[font=\footnotesize] at (8.1,1.3) {\(C\)};
  \node[left,font=\footnotesize] at (-0.12,1.3) {读入 \(A\)};
  \node[right,font=\footnotesize,black!65] at (9.15,1.3) {\([0.5,\,0.65)\)，宽 \(0.15\)};
  \draw[black!35,dashed] (4.5,2.2) -- (0,1.6);
  \draw[black!35,dashed] (7.2,2.2) -- (9,1.6);
  % 第三步
  \fill[black!6] (0,-0.2) rectangle (9,0.4);
  \fill[blue!25] (7.2,-0.2) rectangle (9,0.4);
  \draw[black!50] (0,-0.2) rectangle (9,0.4);
  \draw[black!50] (4.5,-0.2) -- (4.5,0.4);
  \draw[black!50] (7.2,-0.2) -- (7.2,0.4);
  \node[font=\footnotesize] at (2.25,0.1) {\(A\)};
  \node[font=\footnotesize] at (5.85,0.1) {\(B\)};
  \node[font=\footnotesize] at (8.1,0.1) {\(C\)};
  \node[left,font=\footnotesize] at (-0.12,0.1) {读入 \(C\)};
  \node[right,font=\footnotesize,black!65] at (9.15,0.1) {\([0.62,\,0.65)\)，宽 \(0.03\)};
  \draw[black!35,dashed] (0,1.0) -- (0,0.4);
  \draw[black!35,dashed] (4.5,1.0) -- (9,0.4);
  \node[font=\footnotesize,black!70,align=center] at (4.5,-0.95)
    {区间宽度 \(=0.3\times0.5\times0.2=P(BAC)\)};
\end{tikzpicture}

\emph{每读一个符号，当前区间按同一套概率比例重切一次，取对应的子段（灰底为放大后的当前区间）。三步之后区间宽 \(0.03\)，恰是这条序列的概率。写下一个落在其中的数所需的比特数，约为 \(-\log_2 0.03\approx5.1\)。}
\end{center}

图上那个巧合不是巧合。每一步宽度乘以当前符号的条件概率，所以任何序列都满足

\[
\abs{I(x^n)}=\prod_{i=1}^{n}P(x_i\given x^{i-1})=P(x^n).
\]

序列的概率就是它的区间宽度。上一节那条``码长等于 \(-\log_2\) 概率''的对应，在这里以连续的形式重新出现：区间越窄，定位它需要的比特越多，而且不必是整数。

顺便注意公式里写的是条件概率。算术编码从不要求符号独立，它只要求编码器和解码器在每一步用完全相同的条件分布。模型可以随上下文任意变化，这一点后面会很重要。

\subsection{把区间兑换成比特，只多付两位}

区间是实数对象，发出去的必须是有限比特。收端要的不是``区间里的某个实数''——实数没有宽度，说不清楚——而是一段二进制前缀，它所代表的整个二进小区间完整地落在目标区间之内。长度为 \(k\) 的前缀对应一个宽 \(2^{-k}\) 的格子；比如 \(101000\) 对应 \([0.625,\,0.640625)\)，整个落在 \([0.62,\,0.65)\) 里，于是这六位就唯一标识了 \(BAC\)。

需要多少位？取 \(k=\lceil-\log_2 w\rceil+1\)，其中 \(w\) 是区间宽度，此时 \(2^{-k}\le w/2\)。宽度 \(2^{-k}\) 的格子等距铺满 \([0,1)\)，像一把刻度均匀的尺子。若目标区间连一个完整格子都容不下，它就只能从两个相邻格子各切一小片，两片合起来不足 \(2\cdot 2^{-k}\le w\)，可它自己的宽度就是 \(w\)——矛盾。所以至少有一个格子整个躺在里面。

由此得到总码长的上界

\[
\text{码长}\le -\log_2 P(x^n)+2 .
\]

除以 \(n\)，每符号只比理想值多 \(2/n\) 比特，序列一长就趋于零。逐符号编码那笔一比特一份的零头，被压成了整条消息共付的两个比特。

\subsection{压缩率是模型质量的读数}

模型的概率若和真实分布有偏差，算术编码不会报错，只会安静地把文件写长。对独立同分布的数据，每符号的期望码长正是交叉熵

\[
\frac1n\,\E_P\bigl[-\log_2 Q(X^n)\bigr]=H(P,Q)=H(P)+\KL(P\|Q).
\]

拿刚才那个 \(P=(0.5,0.3,0.2)\) 试三个模型，把一千个符号的文件大小画出来。

\begin{center}
\begin{tikzpicture}[x=1cm,y=1cm,font=\small,>=stealth]
  \fill[blue!35] (0,1.8) rectangle (7.425,2.4);
  \fill[blue!35] (0,0.9) rectangle (7.425,1.5);
  \fill[red!35] (7.425,0.9) rectangle (7.925,1.5);
  \fill[blue!35] (0,0.0) rectangle (7.425,0.6);
  \fill[red!35] (7.425,0.0) rectangle (9.41,0.6);
  \draw[black!50] (0,1.8) rectangle (7.425,2.4);
  \draw[black!50] (0,0.9) rectangle (7.925,1.5);
  \draw[black!50] (0,0.0) rectangle (9.41,0.6);
  \draw[red!70,dashed] (7.425,-0.4) -- (7.425,2.8);
  \node[above,font=\footnotesize,red!70] at (7.425,2.82) {\(H(P)=1.485\) 比特/符号};
  \node[left,font=\footnotesize] at (-0.1,2.1) {模型即真相 \(P\)};
  \node[left,font=\footnotesize] at (-0.1,1.2) {均匀模型 \(Q_1\)};
  \node[left,font=\footnotesize] at (-0.1,0.3) {偏得较远的 \(Q_2\)};
  \node[right,font=\footnotesize,black!70] at (7.55,2.1) {\(1485\) 比特};
  \node[right,font=\footnotesize,black!70] at (8.05,1.2) {\(1585\)，多出 \(\KL=0.100\)};
  \node[right,font=\footnotesize,black!70] at (9.55,0.3) {\(1882\)，多出 \(\KL=0.397\)};
\end{tikzpicture}

\emph{同一段一千符号的数据，换一个概率模型，压出来的文件就变长；红色那截的长度分毫不差地等于 \(\KL(P\|Q)\)。压缩率因此是模型质量的直接读数，而不是某种间接指标。}
\end{center}

这张图解释了现代压缩系统为什么把``概率模型''和``熵编码器''分开写。模型负责每一步给出条件分布，它猜得准不准决定区间收缩得快不快；算术编码只负责把概率兑换成比特，除了那两个常数位不产生额外损耗。压缩得好不好，几乎全部归因于模型。

反过来读同样成立，而且是当下更常用的读法：既然码率等于交叉熵，那么把一个语言模型拿去压一段文本，得到的每字符比特数就是它在这段文本上的交叉熵。困惑度是同一个数的指数形式——以比特计的交叉熵为 \(L\) 时，困惑度为 \(2^{L}\)。报告``每字节 \(0.9\) 比特''和报告``困惑度若干''说的是一回事，见\hyperref[sec:07-Information-Theory/01-Information-and-Entropy/02]{``Shannon 熵与二元熵''}。用压缩率给模型排名因此是有理论依据的，但要记住它只测分布拟合的质量：一个把训练语料背下来的模型压缩率极好，别的什么也说明不了。

模型里有一条不能踩的红线：给真实出现的符号分配零概率。该符号的子区间宽度为零，整条序列的区间随之归零，没有任何二进制前缀能落进去，编码直接中断。这就是各种平滑与逃逸机制存在的原因，它们保护的不是精度，是可行性。

\subsection{解码是同一套动作倒着做}

解码器拿着相同的模型，看收到的数落在初始三段的哪一段，读出第一个符号；再把该段拉伸回 \([0,1)\)，看落在哪一段，读出第二个符号。编码在做的每一步收缩，解码都做一次对应的放大。

正因为完全对称，两端必须一模一样。任何一步的概率取值、舍入方式或状态更新有一丝不同，之后所有区间边界都会错位，而且不会立刻暴露——前几个符号往往还能解对，几十步之后全盘皆错，中间没有任何报警。算术编码实现比理论推导多出来的复杂度，基本都花在维持这份同步上。

真实代码不用无限精度实数，而是两个定宽整数寄存器记住区间的两端。当两端的高位变得一致，这些高位就不再影响后续计算，可以立即输出并把区间左移放大，这叫重归一化。若区间长期卡在中点附近、高位迟迟不统一，低位精度会被逐步抽干，需要提前扩区并记一笔欠账，等高位稳定后再补正。此外必须约定终止方式——传消息长度，或在字母表里加一个专用的结束符号——否则解码器不知道什么时候停下，而最后一个符号常常不产生任何输出比特，这是早期实现里最隐蔽的一类缺陷。概率表也要量化成两端一致的整数累积频数，浮点原型不能直接当跨平台格式，因为不同架构的舍入差异足以让区间偏移。

范围编码是同一思想的另一种整数化路径，工程选择不同，主线不变：序列的概率就是区间的宽度。

到这里，上界的那一侧已经交代完了——熵加上可以忽略的常数总能达到。下界那一侧目前还只有一个逐符号的证明。下一节换一个视角：把长度为 \(n\) 的序列当成一个整体来数，看看``实际会出现的序列''到底有多少个。数出来的那个指数，正是熵。
